\chapter{System Design and Technical Implementation}
\label{ch:system_design_implementation}

This chapter details the architectural design and implementation of SENTINEL: the two-tier execution pipeline, Tier 1 wire-speed filtration and Machine Learning Gate, the Tier 2 agentic reasoning engine and Threat Memory, Dual-RAG knowledge retrieval, adversarial security controls with cryptographic audit sealing, and the local deployment architecture.

\section{SENTINEL Architecture Overview and SOC Stack Integration}
\label{sec:architecture_overview}

SENTINEL enforces a Cognitive Two-Tier Architecture whose core design principle is strict decoupling of work by cost: Tier 1 executes stateless statistical and machine learning algorithms resolving the majority of traffic at wire-speed, while Tier 2 hosts a stateful language model agent performing deep contextual reasoning for complex anomalous events.

Figure~\ref{fig:sentinel_flow} illustrates the execution flow. Raw telemetry from perimeter sources (NGFW, IPS, EDR) is ingested into a Redis Stream queue. On arrival, Tier 1 computes online $Z$-score metrics using Welford's running statistics and classifies feature interactions through the LightGBM Machine Learning Gate. Flows conforming to baseline distributions, or matching query templates already held in the Tier 1.75 Semantic Cache, are adjudicated sub-millisecond without touching the GPU.

\begin{figure}[H]
    \centering
    \includegraphics[width=0.66\textwidth]{images/sentinel_flow.png}
    \caption{Operational execution workflow within the SENTINEL pipeline. (Source: Author)}
    \label{fig:sentinel_flow}
\end{figure}

One design decision must be stated explicitly: among Tier 1's action vocabulary, \textbf{only \texttt{ESCALATE} proceeds to Tier 2}. The adjudications \texttt{BLOCK\_IP}, \texttt{ALERT}, \texttt{AWAIT\_HITL}, and \texttt{DROP} all \textbf{terminate at Tier 1}. This is the basis for the offload definition used in Chapter 4: an event counts as offloaded when it generates no language model invocation, regardless of whether the final adjudication was to block or to pass.

At Tier 2, the LangGraph agent triggers a reasoning cycle on the locally hosted Foundation-Sec-8B-Instruct model, integrating historical state from Threat Memory and retrieved evidence from Dual-RAG. Final decisions are governed by a deterministic guardrail bounded to three states, and every record is sealed with an HMAC-SHA256 hash chain. Within an enterprise SOC, SENTINEL therefore acts as cognitive middleware between the SIEM collector and downstream firewall/SOAR execution platforms: it absorbs stream noise at Tier 1 and supplies grounded, cited justifications for escalated events, while guaranteeing data sovereignty through fully air-gapped execution on a single 16\,GB VRAM GPU.

\section{Tier 1: Wire-speed Filter and Machine Learning Gate}
\label{sec:tier1_mlgate}

Tier 1 is the stateless perimeter layer, combining the Welford $\mathcal{O}(1)$ statistical baseline filter (Equation~\ref{eq:welford}) with a LightGBM Machine Learning Gate to guarantee millisecond latency before any GPU invocation.

The statistical filter computes $Z = (x - \mu)/\sigma$ in real time across flow attributes (packet rates, flow volumes, connection durations), with $\mu$ and $\sigma$ updated via the recurrences of Section~\ref{sec:welford} without retaining raw series~\cite{welford1962}; traffic conforming to the baseline is shed immediately. Alongside this statistical branch, Tier 1 operates an application signature set covering \textbf{29 web attack families} with multi-round normalisation in the OWASP CRS style---decoding URL and HTML entities before matching---to catch payloads disguised by encoding. For repetitive volumetric behaviour, the Tier 1.75 Semantic Cache hashes normalised query strings in RAM and reuses existing adjudications.

A LightGBM Gradient Boosted Decision Tree model follows, evaluating non-linear feature interactions. It uses leaf-wise growth and is trained on 949,535 NetFlow samples (664,674 training, 189,907 hold-out). To neutralise evasion, the Gate clips feature $Z$-scores within $\pm 8\sigma$ and imputes missing values with baseline means; if distorted feature proportions exceed 30\%, it \textbf{abstains} and escalates. Abstention is deliberate: a confident wrong answer is more dangerous than an admission of ignorance, because it causes the system to act autonomously with no one reviewing. Gate adjudications map to actions through a unified \textbf{four-band confidence policy} ($0.85$ / $0.65$ / $0.40$ for \texttt{BLOCK\_IP} / \texttt{ESCALATE} / \texttt{ALERT} / \texttt{LOG}), so that confidence \textit{drives} the action rather than merely accompanying it, and autonomous blocking is isolated to the top band, which additionally requires Tier 1 rule concurrence.

\section{Tier 2: Stateful LangGraph Agent and Threat Memory}
\label{sec:tier2_agent_fsm}

Tier 2 performs stateful cognitive reasoning, executing a GGUF Q4\_K\_M quantized Foundation-Sec-8B-Instruct model~\cite{foundation_sec_2024} via CUDA-compiled llama.cpp on a single 16\,GB VRAM GPU, cutting required VRAM from ~18\,GB (FP16) to 7--8\,GB with negligible accuracy loss.

Reasoning is orchestrated by a LangGraph state-graph~\cite{langgraph2024} across five nodes: guardrail inspection, context enrichment (\texttt{rag\_\allowbreak context}), triage (\texttt{llm\_\allowbreak triage}), technique mapping (\texttt{attack\_\allowbreak mapper}), and action execution or hand-off to \texttt{human\_\allowbreak in\_\allowbreak loop}. Technique mapping requires explicit verification from domain knowledge; absent an unambiguous match, the graph routes to analyst review (\texttt{AWAIT\_HITL}) rather than emitting a speculative action. The human-review path \textbf{is a designed output, not a failure}: the system triages the work requiring a human rather than eliminating the human. At \texttt{llm\_\allowbreak triage} the system performs \textbf{label stripping}, removing every internal classification field from the prompt---a precondition for the validity of every measurement in Chapter 4, since a ground-truth label reaching the context would let the model copy the answer rather than reason to it.

To correlate traces over time, Tier 2 maintains Threat Memory on an index-optimized SQLite database configured with \texttt{synchronous=\allowbreak NORMAL} and rollback-journal rather than WAL, preventing cross-UID file lock crashes under Docker. It tracks IP reputation (\texttt{ip\_\allowbreak reputation}), internal entities (\texttt{known\_\allowbreak entities}), risk indicators (\texttt{apt\_\allowbreak indicators}), and threat events (\texttt{threat\_\allowbreak events}). Reputation score $S$ decays over silent period $t$:
\begin{equation}
    S_{\text{new}} = S_{\text{old}} \times (1 - \lambda)^t
\end{equation}
where $\lambda$ is the decay rate. On reaching the maximum ($S = 100$), Tier 1 enforces block-on-sight for subsequent connections without re-triggering Tier 2---closing an automated feedback loop that converts a one-time reasoning cost into a reusable adjudication. Internal infrastructure addresses are whitelisted and immune to blocking, preventing self-DoS.

\section{Dual-RAG Hybrid Knowledge Retrieval}
\label{sec:dual_rag}

To limit hallucination and ground security decisions, Tier 2 incorporates a Dual-RAG retrieval system over 433 STIX 2.1-structured MITRE ATT\&CK entries~\cite{strom2018mitre} and NIST SP 800-61r2 guidance~\cite{nist80061}. Retrieval runs two complementary mechanisms: dense semantic search via FAISS~\cite{faiss2019} with a Sentence-BERT embedding model~\cite{reimers2019sentence}, identifying techniques sharing semantic intent despite keyphrase variation; and sparse keyword matching via BM25~\cite{bm25_2009}, ensuring exact retrieval of identifiers such as CVE codes, Registry keys, or service ports. Candidates are consolidated using RRF (Equation~\ref{eq:rrf}) with $k=60$ and injected into the prompt at \texttt{rag\_\allowbreak context}.

On top of retrieval the system applies an \textbf{evidence-anchoring guardrail}: every technique identifier emitted by the model is checked back against that batch's retrieved documents, and identifiers absent from them are forced to \texttt{N/A} and routed to \texttt{AWAIT\_HITL}. This is a \textbf{verifiable architectural constraint} rather than a dependency on how prone a given model is to hallucinate---and it is precisely what separates a knowledge-grounded system from a wrapper around a language model.

\section{Adversarial AI Security and Cryptographic Audit Sealing}
\label{sec:ai_security_hmac}

Log-substrate prompt injection~\cite{promptinjection2023} embeds control instructions inside raw telemetry fields (User-Agent headers, URI parameters, Syslog text) to hijack model behaviour. SENTINEL implements Delimited Data Encapsulation with a cryptographic nonce regenerated per processing batch~\cite{owasp_llm2025}: raw telemetry is enclosed within nonce-bearing delimiters, forcing the model to treat the enclosed content as passive data. Per-batch generation---rather than a fixed delimiter---is essential, because an invariant delimiter can be inferred from source code or repeated observation, at which point an attacker can close the encapsulation and escape it. In parallel, a deterministic guardrail enforces Pydantic schema validation (\texttt{LLMDecision}); on formatting errors, regular expression salvaging recovers essential fields or defaults to the safe \texttt{AWAIT\_HITL} state.

For auditing, each adjudication record $i$ is chained using HMAC-SHA256~\cite{hmac1997}:
\begin{equation}
    H_i = \text{HMAC-SHA256}(D_i \parallel H_{i-1}, K)
\end{equation}
where $D_i$ is the current record, $H_{i-1}$ the preceding hash, and $K$ a local secret key loaded from the environment rather than hardcoded, since a key exposed in source allows the entire chain to be forged from scratch. Chaining ensures any alteration or deletion within the chain breaks continuity and is flagged immediately. As stated in Section~2.4, this guarantees \textbf{integrity and tamper-evidence}; because the key is symmetric, it does not constitute non-repudiation in the cryptographic sense.

\section{SOC Dashboard and On-Premises Deployment Architecture}
\label{sec:dashboard_deployment}

The SOC Analyst Dashboard is built with Streamlit, connecting directly to Redis Stream and Threat Memory to provide real-time visibility: the per-tier offload funnel, multi-stage intrusion timelines, HMAC chain integrity status, and the Human-in-the-Loop review panel. All services are containerized with Docker Compose; \texttt{llama.\allowbreak cpp} is compiled with CUDA for the RTX 4060 Ti; and \texttt{threat\_memory.db} with the FAISS index are mounted on persistent volumes so reputation scores and audit records survive restarts.

Deployment adheres to Zero Trust mandates~\cite{nist80207}: password strings in \texttt{REDIS\_URL} are redacted before logging. To manage stream surges, backpressure is driven by \textbf{Redis consumer group lag} rather than queue length. This choice is mandatory: a Redis Stream's length only grows over time even when every record has been fully processed, so using it as the throttling signal would strangle the producer permanently; only consumer group lag reflects the work actually outstanding.
